% ===================================================================
%  TAULA N° 11 — La Vila e sos monuments. — L'Incendi.
%  Traduction 2026 d'apres texto/io, controlee sur texto/fr et
%  texto/ca. Consigne : voir texto/oc/10-taula-01.tex.
% ===================================================================

%%K t11-apar-1 apar
\VUsaut{2.0mm}
\VUcentre{13.2pt}{90}{TRESENA SÈRIA}
\VUsaut{3.0mm}
\VUfilet{16mm}
\VUsaut{5.6mm}
\VUtitre{15.0pt}{60}{TAULA N° 11}
\VUsaut{1.4mm}
\VUpk{11.4pt}{40}{La Vila e sos monuments. --- L'Incendi.}
\VUsaut{2.2mm}

%%K t11-tit-1 sub
\VUpk{10.2pt}{40}{Los faubors.}
\VUsaut{3.0mm}

%%K t11-c1-01-1 p
1. --- Demòri dins una granda vila situada a cent quilomètres de l'ocean e que,
çaquelà, es un \VUgras{pòrt} \textsuperscript{(1)} de primièr òrdre. Lo riu que la
banha a quatre cents mètres de larg e forma una rada fòrça bèla.

%%K t11-c1-02-1 p
2. --- Tota la partida de la vila qu'es contra los \VUgras{cais}
\textsuperscript{(2)} a un aspècte tot maritim e industrial e forma un
\VUgras{faubor} \textsuperscript{(3)} abitat pas que per de mariniers e
d'obrièrs. Aqueles trabalhan dins las \VUgras{fabricas} \textsuperscript{(4)} e
dins la \VUgras{fabrica de gas} \textsuperscript{(5)}, que sa \VUgras{chiminèa}
\textsuperscript{(6)} nauta e son \VUgras{gasomètre} se veson de fòrça luènh.

%%K t11-c1-03-1 p
3. --- Dins l'Edat Mejana la vila èra fortificada, e se tròban encara qualques
rèstas de sas muralhas, en particular la \VUgras{pòrta} \textsuperscript{(8)} del
vièlh \VUgras{castèl fòrt} \textsuperscript{(7)}, flanquejat per sas doas
\VUgras{torres} \textsuperscript{(9)}, qu'uèi servisson de \VUgras{preson}.

%%K t11-c1-04-1 p
4. --- Dins tot aqueste \VUgras{quartièr}, d'aspècte fòrça ancian, un sol
bastiment es relativament modèrne: la \VUgras{doana} \textsuperscript{(10)}. Es
entre lo cai e lo pichon \VUgras{carrièr} \textsuperscript{(11)} que mena a la
vièlha \VUgras{comuna} \textsuperscript{(12)}. Aqueste darrièr monument se
reconeis aisidament pel gigantesc \VUgras{cloquièr} \textsuperscript{(13)} que
domina tota la vila vièlha.

%%K t11-c1-05-1 p
5. --- De l'autre costat del carrièr \VUgras{s'auça} un bastiment bastit dins lo
meteis estil que la doana: es la \VUgras{bòrsa} \textsuperscript{(14)}. Una de sas
façadas dona sus una placeta ont es la \VUgras{prefectura} \textsuperscript{(15)}.
Nòstra vila, en efièch, sens èsser una capitala, es un important cap de
departament.

%%K t11-tit-2 sub
\VUsaut{5.0mm}
\VUpk{10.2pt}{40}{La Partida nauta de la Vila.}
\VUsaut{3.4mm}

%%K t11-c1-06-1 p
6. --- La garnison, pro nombrosa, s'alòtja dins de vastas \VUgras{casèrnas}
\textsuperscript{(16)} fòrça salubras; s'espandisson fins a la reissa del
\VUgras{pargue municipal} \textsuperscript{(17)}. Lo \VUgras{baloard}
\textsuperscript{(19)} que mena a aquel pargue passa darrièr la \VUgras{caissa
d'estalvis} \textsuperscript{(18)} e desemboca sus la plaça de la
\VUgras{catedrala} \textsuperscript{(20)}.

%%K t11-c1-07-1 p
7. --- Totes aqueles monuments s'escalonan en amfiteatre sus un truc ont es tanben
nòstre larg \VUgras{institut} \textsuperscript{(21)}.

%%K t11-c1-07-2 p
Se se davala per una via un pauc en pendís que se junta amb lo Grand Carrièr,
s'arriba al \VUgras{palais de justícia} \textsuperscript{(22)}, davant lo qual
s'espandís un agradiu \VUgras{jardin public} \textsuperscript{(26)}. Aquesta
\VUgras{plaça jardinada} es pas fòrça granda, mas pausa al mitan de la vila un
oasi de verdura e de frescor. Per aquò la frequentan fòrça los abitants, que lor
agrada de venir s'assetar jos lo tendal del \VUgras{cafè} \textsuperscript{(25)}
per escotar la musica militara que tòca lo dijòus e lo dimenge dins lo
\VUgras{quiòsc} \textsuperscript{(27)} pausat entre las gèspas e los massisses.

%%K t11-c1-08-1 p
8. --- Sus una d'aquelas gèspas s'auça una \VUgras{estatua}
\textsuperscript{(28)} de bronze, monument bastit a la memòria d'un de nòstres
concitadins que rendèt de grands servicis a sa vila natala.

%%K t11-tit-3 sub
\VUsaut{4.0mm}
\VUpk{10.2pt}{40}{Los transpòrts.}
\VUsaut{3.4mm}

%%K t11-c1-09-1 p
9. --- Nòstra vila es pas encara enlumenada amb la lutz electrica; a pas que de
\VUgras{fanals de gas} \textsuperscript{(29)}. Mas la \VUgras{premsa locala} fa
campanha per qu'aquel sistèma d'enlumenatge se melhore, \VUgras{coma} s'es ja
perfeccionat lo \VUgras{sistèma de traccion} dels tramvais. \VUgras{Las} vièlhas
veituras \VUgras{tiradas} per de \VUgras{cavals} son estadas remplaçadas per de
practics \VUgras{tramvais electrics} \textsuperscript{(31)}, que menan rapidament
los viatjaires d'un cap de \VUgras{la vila a l'autre} per un prètz fòrça
\VUgras{barat}.

%%K t11-c1-10-1 p
10. --- Contra lo palais de justícia i a la \VUgras{banca} \textsuperscript{(32)},
ont se descomptan los valors comercials. Los \VUgras{garçons de banca}
\textsuperscript{(33)} van a domicili far los encaissaments.

%%K t11-tit-4 sub
\VUsaut{4.0mm}
\VUpk{10.2pt}{40}{Art e Instruccion.}
\VUsaut{3.4mm}

%%K t11-c1-11-1 p
11. --- Lo bèl bastiment qu'es al costat es lo \VUgras{musèu}. Conten a l'encòp la
\VUgras{bibliotèca} \textsuperscript{(34)} de la vila, de bèlas
\VUgras{colleccions de sciéncias naturalas} \textsuperscript{(35)} (musèu
d'istòria naturala) e una graciosa \VUgras{galariá de pintura}
\textsuperscript{(36)} que constituís una esplendida \VUgras{exposicion}
permanenta.

%%K t11-c1-11-2 p
Es dobèrt al public dos còps per setmana, mas los estrangièrs lo pòdon visitar
quin jorn que siá en donant una estrena al \VUgras{garda} \textsuperscript{(37)}.
Los \VUgras{pintres} \textsuperscript{(38)} i venon trabalhar. Los òm vei
\VUgras{davant} son cavalet, amb la \VUgras{paleta} a la man, a copiar los
tablèus celèbres.

%%K t11-c1-12-1 p
12. --- Sus l'autura, darrièr lo musèu, se destria la \VUgras{cupòla}
\textsuperscript{(40)} de l'\VUgras{espital} \textsuperscript{(39)} e los
bastiments de l'\VUgras{escòla normala} \textsuperscript{(41)} ont se formavan los
mèstres de nòstras escòlas municipalas. Sus la plaça del Teatre i a un
\VUgras{burèu de pòsta} \textsuperscript{(43)} installat dins un \VUgras{ostal
particular} \textsuperscript{(42)}.

%%K t11-tit-5 sub
\VUsaut{5.0mm}
\VUpk{11.4pt}{40}{L'Incendi.}
\VUsaut{3.0mm}

%%K t11-tit-6 sub
\VUpk{10.2pt}{40}{Lo crit de «Al fuòc!».}
\VUsaut{3.4mm}

%%K t11-c2-01-1 p
1. --- Una novèla alarmanta s'escampa per la vila: ven d'esclatar un incendi al
teatre! Quand arribi sus la \VUgras{plaça} \textsuperscript{(96)}, la multitud
s'atropèla ja sul \VUgras{trepador} \textsuperscript{(46)} e sus la
\VUgras{caussada} \textsuperscript{(47)}. Los \VUgras{emplegats de pòsta}
\textsuperscript{(44)} e los \VUgras{factors} \textsuperscript{(45)} sortisson del
burèu. Un \VUgras{cochièr de tramvai} \textsuperscript{(48)} a degut arrestar son
tramvai per daissar passar una \VUgras{ambulància} \textsuperscript{(50)}
municipala qu'arriba amb un \VUgras{cirurgian} \textsuperscript{(52)} e un
\VUgras{infirmièr} \textsuperscript{(51)} per suènhar un \VUgras{ferit}
\textsuperscript{(53)} qu'an portat sus una \VUgras{civièira}
\textsuperscript{(54)} fins a prèp del \VUgras{quiòsc de jornals}
\textsuperscript{(55)}.

%%K t11-c2-02-1 p
2. --- Una companhiá d'infantariá que tornava de l'exercici s'es arrestada per
mantenir l'òrdre amb los \VUgras{gardas municipals} \textsuperscript{(61)} a
caval. Los \VUgras{cavalièrs}, que lors cavals se cabran, o fan melhor que los
\VUgras{soldats} \textsuperscript{(57)} o los \VUgras{gendarmas}
\textsuperscript{(63)} per far recular los \VUgras{curioses}
\textsuperscript{(56)}. Autrament, aqueles darrièrs son fòrça ocupats. Un d'eles
indica als \VUgras{agents de polícia} \textsuperscript{(64)} ont cal menar una
autra \VUgras{victima} \textsuperscript{(65)}, un valent pompièr qu'es tombat
d'una escala.

%%K t11-c2-03-1 p
3. --- L'\VUgras{oficièr} \textsuperscript{(66)} precisa l'endrech ont se deu
pausar la \VUgras{bomba de vapor} \textsuperscript{(67)} qu'arriba. En esperant
aquela poderosa maquina, a fach montar una \VUgras{bomba a man}
\textsuperscript{(68)}, que son long \VUgras{tuble}
\textsuperscript{(69)} geta de rajòls d'aiga sul fuòc. Sièis pompièrs la manejan
mentre que lor companh, que ten la \VUgras{lança} \textsuperscript{(70)},
dirigís lo \VUgras{rajòl} \textsuperscript{(71)} al centre de l'incendi.

%%K t11-c2-04-1 p
4. --- De l'autre costat del teatre s'es apiejada la granda \VUgras{escala}
\textsuperscript{(72)}. Permet als pompièrs de s'apropchar de las flamas que
brotan entre de revolums de \VUgras{fum} \textsuperscript{(75)} negre.

%%K t11-tit-7 sub
\VUsaut{5.0mm}
\VUpk{10.2pt}{40}{Actes de valor.}
\VUsaut{3.4mm}

%%K t11-c2-05-1 p
5. --- L'\VUgras{incendi} \textsuperscript{(74)} nasquèt dins lo vestibul del
\VUgras{teatre} \textsuperscript{(73)}, mentre que s'assajava l'\VUgras{opèra} que
sa primièra \VUgras{representacion} deviá se tenir deman. Per bonur i aviá pas
qu'un pauc d'\VUgras{espectators} \textsuperscript{(76)} dins la sala. Los paures
se son refugiats sul \VUgras{balcon} \textsuperscript{(77)}, ont de valents
\VUgras{pompièrs} \textsuperscript{(86)} acorron per los socórrer amb una escala e
de còrdas.

%%K t11-c2-06-1 p
6. --- Jos lo \VUgras{peristil} \textsuperscript{(78)}, ont l'\VUgras{afiche}
\textsuperscript{(79)} anóncia la representacion, se vei fugir los
\VUgras{musicaires} \textsuperscript{(81)} de l'\VUgras{orquèstra}
\textsuperscript{(80)} en emportar lors instruments: \VUgras{violon}
\textsuperscript{(81)}, \VUgras{flaüta} \textsuperscript{(82)},
\VUgras{violoncèl} \textsuperscript{(83)}, \VUgras{trombòn}
\textsuperscript{(84)}, \VUgras{tambor} \textsuperscript{(85)}, etc. S'ensaja de
sauvar una partida del \VUgras{moblament} \textsuperscript{(87)}.

%%K t11-c2-07-1 p
7. --- Los \VUgras{actors} \textsuperscript{(89)} e las \VUgras{actrises}
\textsuperscript{(88)}, \VUgras{cantaires} eles e elas, an degut fugir amb lor
\VUgras{vestit} \textsuperscript{(90)} de teatre. Lo tenòr explica a un
\VUgras{jornalista} \textsuperscript{(92)} cossí es arribat. \VUgras{Lo reportèr}
acorreguèt del primièr moment amb las autoritats: lo \VUgras{prefècte}
\textsuperscript{(91)}, lo \VUgras{cònsol} \textsuperscript{(93)}, lo
\VUgras{comissari de polícia} \textsuperscript{(94)} e un \VUgras{jutge}
\textsuperscript{(95)}, que dobriràn una enquèsta per determinar las
responsabilitats.
\VUsaut{6.0mm}
\VUfilet{22mm}
